\section{Complete Multi-Signal Research Iteration Index}

All twenty iteration notes below belong to \texttt{feature/volume-anomaly} at
\texttt{4279d84}. They are included to preserve positive, negative, numerical,
transfer, and operational evidence. No iteration is silently incorporated into
the reference detector, and these records were read rather than rerun for this
report. Later changes must not inherit the final-test claim of iteration 005.

\begin{longtable}{>{\raggedright\arraybackslash}p{0.45\textwidth}>{\raggedright\arraybackslash}p{0.47\textwidth}}
\caption{Research iteration records and their decisions.}\label{tab:research-iterations}\\
\toprule\textbf{Record} & \textbf{Role in the project} \\
\midrule\endfirsthead
\toprule\textbf{Record} & \textbf{Role in the project} \\
\midrule\endhead
\path{reports/UNSW-Chromecast/iteration-001-baseline.md} & Records the four-category DM baseline and the attacks it misses. \\
\path{reports/UNSW-Chromecast/iteration-002-volume-hybrid.md} & Adds a high packet-count tail; recovers a flood without another development false alert. \\
\path{reports/UNSW-Chromecast/iteration-003-behavior-tails.md} & Adds byte, SYN, and direction behaviour checks; improves attack coverage. \\
\path{reports/UNSW-Chromecast/iteration-004-arp-tail.md} & Adds an ARP-count tail for low-rate ARP attacks. \\
\path{reports/UNSW-Chromecast/iteration-005-locked-final.md} & Freezes an earlier detector and records its one-time final evaluation. \\
\path{reports/D-LinkDayCam5/iteration-006-cross-device-normal.md} & Tests transfer on unlabelled D-Link traffic; alerts are not measured FPR. \\
\path{reports/UNSW-Chromecast/iteration-007-silent-window-abstention.md} & Avoids treating an undefined direction in silent windows as an anomaly. \\
\path{reports/UNSW-Chromecast/iteration-008-temporal-incident-evaluation.md} & Adds offline incident/delay metrics; rejects noisy EWMA and port-churn features. \\
\path{reports/UNSW-Chromecast/iteration-009-edge-benchmark.md} & Measures scoring latency, state size, and process memory on the development host. \\
\path{reports/UNSW-Chromecast/iteration-010-guarded-adaptation.md} & Builds an anchored shadow candidate; eligibility does not automatically promote it. \\
\path{reports/UNSW-Chromecast/iteration-011-native-lm-backend.md} & Uses native Languasco--Migliardi likelihood with numerical/decision agreement. \\
\path{reports/D-LinkDayCam5/iteration-012-long-horizon-drift.md} & Exposes persistent directional drift over 1,218 later D-Link windows. \\
\path{reports/UNSW-Chromecast/iteration-013-two-block-direction-dm.md} & Replaces raw direction fraction with a two-category direction DM model. \\
\path{reports/UNSW-Chromecast/iteration-014-jackknife-robustness.md} & Checks omitted fit/calibration days; preserves recall with limited false-alert variation. \\
\path{reports/UNSW-AmazonEcho/iteration-015-third-device-transfer.md} & Applies the frozen architecture to Amazon Echo without device-specific tuning. \\
\path{reports/iteration-016-calibrated-tail-fusion.md} & Calibrates feature score ranks; keeps decisions while increasing retained threshold size. \\
\path{reports/iteration-017-temporal-byte-change-rejection.md} & Rejects adjacent-window byte change because of false alerts and transfer cost. \\
\path{reports/iteration-018-levenberg-marquardt-audit.md} & Rejects a true Levenberg--Marquardt optimiser after a valid synthetic failure. \\
\path{reports/iteration-019-mean-packet-size-rejection.md} & Rejects mean-packet-size tails after cross-device false-positive growth. \\
\path{reports/iteration-020-corroboration-rule-rejection.md} & Rejects corroboration rules that suppress valid single-signal attacks. \\
\bottomrule
\end{longtable}
